\section{Consequences and scope}


\paragraph{The coefficient reads the spine.}
For \(w=h=0\), a bipartition \(A,B\) of the spine and the change \(y_i=t_i\) on \(A\), \(y_j=1-t_j\) on \(B\) turn \(\mathcal P\) into the order polytope for the two-level poset with \(i<j\) exactly on edges \(ij\) from \(A\) to \(B\). Its volume is the number of linear extensions divided by \(m!\): ordering the \(m\) coordinates triangulates the cube into simplices of equal volume, and exactly the linear extensions satisfy the inequalities \cite{Stanley}. The star gives \(1/m\) by integrating \((1-t)^{m-1}\) over its centre. The path gives the fraction of alternating permutations, in particular \(1/2\) for \(m=2\) and \(5/24\) for \(m=4\). The \(m=4\) star and path manifolds therefore have unequal WRT coefficients despite \(H_1=(\ZZ/4\ZZ)^8\). The \(m=2\) untwisted graph of Figure~\ref{fig:plumb} has coefficient \((-1)^r/(2\pi^4)\) in \(r^4(\log r)^2\). With a single nonzero spine twist its volume is \(3/8\); with a single nonzero pendant twist it is \(1/8\). Neither example is a selected observable of a different manifold: the changed twist is a change of the physical framing in surgery.

\paragraph{The logarithm-free formulation does not cover this family.}
For the untwisted manifolds, the leading power \(r^{2m}(\log r)^m\) has nonzero coefficient and phase \((-1)^r\). Suppose an ordinary Poincar\'e expansion were a finite sum of fixed unit phases with Laurent--Puiseux powers and with logarithmic degrees bounded by \(b_0<m\). The expansion cannot be identically zero, as the invariant grows. Combine coincident phases, and take the largest nonzero pair \((d,b)\) in the lexicographic order of \(r^d(\log r)^b\). After division by its scale the putative leading phase sum is bounded. If \((d,b)<(2m,m)\), the actual quotient is unbounded. If \((d,b)>(2m,m)\), the quotient tends to zero, which a nonzero finite sum of distinct phases cannot do: its Ces\`aro mean square is the sum of the squared moduli of its coefficients. The equal pair is excluded by \(b_0<m\). The argument also applies on either parity after merging phase classes on that progression. As \(m\) is arbitrary, no uniform degree bound holds across closed manifolds in expansions of this form.

In particular this excludes the literal all-closed Laurent--Puiseux formulation in Andersen, Han, Li, Misteg\aa rd, Sauzin and Sun, Conjecture~1.1 of arXiv:2510.10678v1 \cite{AHLMS}, in its stated \(\tau_r(S^3)=1\) normalisation. It does not concern their Theorem~1.1 on Seifert \emph{integral} homology spheres. Log powers are already allowed in certain analytic expansions, such as Andersen--Petersen's mapping-torus results \cite[Theorems~1.3 and 3.4]{AP}. We do not assert a full log-Puiseux expansion, determine the first surviving term when \(\vol(\mathcal P)=0\), or attribute the asymptotics to a specific flat-connection stationary action. No priority claim is made.


\begin{corollary}[Obstruction to definite tree plumbing]\label{cor:definite}
If \(\vol(\mathcal P)>0\), neither \(M_{\mathcal T;w,h}\) nor its orientation reverse admits a negative-definite weighted-tree plumbing presentation in the class of Murakami~\cite[Theorem~1.2 and \S\S\,9.1--9.3]{Murakami}. Consequently neither admits a positive-definite presentation of the same tree-plumbing type.
\end{corollary}
\begin{proof}
Murakami's Theorem~1.2 gives a finite sum of fixed phases with Laurent half-power coefficients and no logarithms for negative-definite plumbed manifolds in that class. His \S\,9.2 uses the same \(S^3=1\) normalisation and the root convention of the Gukov--Pei--Putrov--Vafa surgery formula, with his \(k\) our shifted index \(r\). On the other hand, Theorem~\ref{thm:corona} gives \(|\tau_r(M)|\sim c\,r^{2m}(\log r)^m\) for \(c=\pi^{-2m}\vol(\mathcal P)>0\). The preceding finite-phase argument excludes Murakami's expansion, irrespective of the unit phase \(\Phi_r\). Orientation reversal complex conjugates the unitary invariant and therefore leaves this obstruction unchanged. Mirroring a positive-definite tree surgery gives a negative-definite presentation of the reverse orientation: reverse the framings and linking signs, then reorient the components on one colour of the bipartite tree if needed to restore positive Hopf clasps.
\end{proof}
This is an obstruction to the existence of a presentation in that plumbing class, not just a statement about the inertia of the displayed surgery matrix. It makes no claim about arbitrary definite four-manifolds that might bound the three-manifold.

The tree theorem depends on the paired opposite order-two caps and the chosen closed tree plumbing. These ingredients preserve one explicitly computable finite transfer; other cap families, arbitrary torus graphs and signed possibly vanishing leading coefficients require separate estimates. In particular we do not deduce a full expansion of a general graph manifold or attribute every term to a flat-connection stratum. The \(S,T\) matrices are shared with marked torus sewing experiments, but no measured open boundary price or cut observable enters the present closed invariant.
